\section{Complete stationary-three reference C and an unbounded binary-dense actual cone}
Let $S=z+z^2+z^3$, $T=z+z^2$, $B_r=\sum_{j<r}S^j$, and
\[
 R^{(r)}=B_r+\frac{S^r}{1-T},\quad W=3r,\quad A=r+3.
\]
This is the infinite-BINARY reference with r final ternary parts,
not the infinite-initial-ones derivative used in the preceding pure
ternary theorem. Put $\phi=1+q$, $q^2+q=1$, $q>0$, and
$H^{(r)}_n=R^{(r)}_n-\phi R^{(r)}_{n-1}$.

\begin{theorem}
For EVERY r>=0, the reference $R^{(r)}$ has full C at every current
rank and every shift at least three. For r>=6 its ordinary coefficients
are LC at EVERY centre A<=n<=W. This ordinary-LC statement stops
at W: the minor at W+1 is strictly negative for every r>=1.

Consequently EVERY ACTUAL $1^\alpha2^u3^r$ with
\[
                 \alpha\ge1,\qquad u\ge\max(0,2r-3)
\]
has full C at all current ranks and all shifts. The binary and ternary
counts and padding length are all unbounded. General C and original
one-descent coefficient unimodality remain OPEN.
\end{theorem}

\paragraph{A complete golden residual with its alternating tail retained.}
The reference satisfies $R^{(r+1)}=1+S R^{(r)}$ and
\[
 H^{(r)}(z)=\frac{C_r(z)}{1+qz},\quad C_r=1+z^3B_r,\qquad
 H^{(r+1)}=S H^{(r)}+(1-\phi z).                 \tag{1}
\]
The polynomial C_r has degree W, top coefficient one, and next
coefficient r-1 for r>=1. Its finite prefix H_0..H_W is not
identified with a polynomial outside that interval. In fact
\[
 \tau_r=H^{(r)}_W=(2q^2)^r>0,\qquad
 H^{(r)}_{W+j}=(-q)^j\tau_r\quad(j\ge0).         \tag{2}
\]
For example evaluate C_r at -phi using $1-T(-\phi)=0$ and
$S(-\phi)=-2\phi$, or induct directly in (1). The complete top
reads are
\[
 H_W=\tau_r,\quad H_{W-1}=\phi(1-\tau_r),\quad
 H_{W-2}=\phi\{r-1-\phi(1-\tau_r)\}.           \tag{3}
\]
The registered all-reference golden-head theorem proves H_n>0
for EVERY 2<=n<=W. These positive entries, not the later alternating
ones, form the finite tail used below.

\paragraph{Unbounded controls for the unperturbed H3 residual.}
Let $h=(1-S)^{-1}$ and $K_n=h_n-\phi h_{n-1}$. It obeys
the tribonacci recurrence with initial values 1,-q,1-q.
For every n>=6,
\[
 0<K_{n-1},\quad K_n\le2K_{n-1},\quad
 |K_n^2-K_{n-1}K_{n+1}|\le K_{n-1}.             \tag{4}
\]
Here is a full induction rather than a sampled envelope. The
positive starting values at ranks 5,6,7 are
6-7q,11-13q,20-24q. The initial values of
$2K_{n-1}-K_n$, at n=6,7,8, are 1-q,2-2q,3-4q.
Both sequences obey the positive tribonacci recurrence, so positivity
propagates. The minors at centres 6,7,8 are 2-3q,-3+q,1+4q,
bounded in absolute value by the respective preceding K values.
Their exterior-square recurrence is $M_{n+3}=-M_{n+2}-M_{n+1}+M_n$.
The registered signed-majorant argument proves the whole (4).
The new Lean theorem verifies this entire assembly, with the bounds
3/5<=q<=5/8 derived from the quadratic equation.

\paragraph{Both new left entries are paid uniformly.}
We prove, by induction from r=6, that
\[
 (H^{(r)}_n)^2\ge H^{(r)}_{n-1}H^{(r)}_{n+1}
                         \quad(A\le n<W).          \tag{5}
\]
The first three ordinary losses of R relative to h, at A,A+1,A+2,
are 1,r+2,(r^2+5r+10)/2. This follows from the first losses
1,2,5 of h relative to $1/(1-T)$ and multiplication by S^r.
Thus the first three H losses relative to K are
\[
           1,\quad b=r+2-\phi,\quad
           c=(r^2+5r+10)/2-\phi(r+2).             \tag{6}
\]
For r>=7, b>=5 and
\[
                c-4b-1=r(r-7)/2+(2-\phi)(r-2)\ge0.
\]
At the first centre, (4) bounds the actual minor below by
$(b-5)K_{A-1}+1\ge0$. At the second centre its exact expansion
is $M_K-2bK_{A+1}+b^2+cK_A+(K_{A+2}-c)$.
This is at least $(c-4b-1)K_A+b^2+(K_{A+2}-c)\ge0$.
The last term is the ACTUAL positive H_{A+2}, since A+2<=W.
Both entrance payments, including this actual-input guard, have a
single scoped Lean theorem.

\paragraph{All middle reads and both terminal corrections.}
Assume (5) at r>=6 and let W=3r. Truncate the old H to the
positive interval [r+2,W]. By (5) this finite sequence is globally
LC. Convolution by S preserves LC. At new centres n>=r+6 and
n<=W, its entire triple agrees with (1); every old read is within
the positive truncated interval. At centre W+1 the preceding and
current values still agree, while the next value is lowered by
q*tau_r, because the first old tail value is -q*tau_r. Lowering that
successor improves the minor.

Only the last new interior centre W+2 remains. Its complete triple
is, by (2),(3),
\[
 \phi(r-1)-1+2\tau_r,\quad
 \phi(1-\tau_r)+q^2\tau_r,\quad 2q^2\tau_r.
\]
The frozen forward-inverse endpoint identity gives its minor
\[
                 \phi^2-2\{\phi+q(r-1)\}\tau_r.      \tag{7}
\]
The OLD pure-ternary forward-CI theorem, for every r>=6, already
proves $2(\phi+qr)\tau_r\le\phi^2$ from its one fixed r=6
potential and the contracting recurrence. Formula (7) has the
additional nonnegative slack 2q*tau_r. Its complete payment is
also kernel checked. The last positive value at W+3 joins zero
extension and has an automatic endpoint minor. Thus every target
centre is covered: two new left entries, all middle reads, the
improved penultimate read, and the one terminal corner.

The ONLY new finite premise is r=6: the eleven values H8..H18
are positive and the nine minors at centres 9..17 are positive.
They are reconstructed exactly in the quadratic field and checked
in ONE Lean report for the exact root; its narrow rational root
bounds are proved internally. No parameter enumeration is used
as an unbounded induction premise.

\paragraph{From residual LC to ordinary LC and complete reference C.}
The first ordinary F minor at A is nonnegative by the previously
registered first-loss scalar kernel and the unbounded ordinary
tribonacci minor bound. In fact its strong bound is at least R_A.
Use $R_n=\phi R_{n-1}+H_n$ and the exact geometric-cumulative
minor-propagation identity of the all-order cut theorem. Positive H
and (5) propagate ordinary LC from A through W-1. At W, the
identity
$M_R(W)=R_WH_W-R_{W-1}H_{W+1}>0$ pays the endpoint directly.
At W+1, H is negative and its successor positive, so the same
identity proves the claimed strictly negative minor.

The first ordinary loss at A is one. The proved H3 renewal-root
residual bound gives $R_A<\lambda_3R_{A-1}$, while every protected
H3 cumulative ratio is at least lambda3. The finite ordinary-to-G
bridge therefore proves every C shift for A<N<=W+1. For N<A
the whole needed prefix is H3, and at N=A the current derivative
alone is lowered, preserving reference C. Finally the OLD general
infinite-binary reference theorem already gives all C at N>=W-1.
These ranges cover every current, without assuming a first mode
for an infinite growing reference. For r<=5, take arbitrarily long
binary prefixes in the OLD nonbinary-weight-at-most-32 complete-C
family; each finite read agrees with the infinite reference. The
case r=0 is the established binary renewal reference.

\paragraph{Return to actual arms with BOTH counts unbounded.}
Let $F_{\alpha,u}=\sum_{j<u}T^j+T^u[\alpha+1]$ be the binary
arm. Then
\[
 R^{(r)}-F_{1^\alpha2^u3^r}
      =S^rT^u\{(1-T)^{-1}-[\alpha+1]\}\ge0.
\]
For u>=1 the first loss occurs exactly at original rank u+r+2:
the bracket begins at degree two, with coefficient two for alpha=1
and one for alpha>=2. If u>=2r-2, every read through W-1 agrees
with the complete-C reference. At u=2r-3 the first ordinary loss
is exactly at W-1; only its current D is lowered in the C comparison
there, so it is still paid. All earlier currents are unchanged.
The OLD universal ACTUAL C theorem covers every N>=W, every
padding and every shift. For r<=5 the weight-32 theorem covers
all u directly; r=0 is already binary. This proves the actual cone.

\paragraph{Formal scope.}
Four generic reports verify the unbounded K controls, both left
payments, inherited terminal budget, and every two-ended window
read guard. One fixed-base report certifies the complete r=6
quadratic-field entrance. Existing forward-CI potential, LC
convolution, ordinary-to-G bridge, renewal bounds, general reference
late C and universal actual late C are reused and not recounted.
Generating functions, exact source losses and the reference/actual
assembly are written proofs, not end-to-end Lean composition
formalizations. No general original UI conclusion is asserted.
